\documentclass[10pt,a4paper]{amsart}
\usepackage[margin=19mm]{geometry}
\usepackage{amsmath,amssymb,mathtools}
\usepackage[scr=rsfs,cal=euler]{mathalfa}
\usepackage{xfrac}
\usepackage[unicode,hidelinks,bookmarks=false]{hyperref}
\newcommand{\ie}{i.e.\ }
\newcommand{\cf}{cf.\ }
\newcommand{\resp}{respectively\ }
\newcommand{\Subsection}{\subsection}
\providecommand{\nobreakdash}{\nobreak\hbox{-}}
\setlength{\emergencystretch}{2em}
\multlinegap0pt
\usepackage{amssymb}
\usepackage{bm}
\usepackage{mathtools}
\usepackage[shortlabels]{enumitem}
\usepackage{tikz}
\newtheorem{lemma}{Lemma}
\title{Euler-type recurrence relations for partition functions with congruence conditions}
\author{Wissam Raji, Hasan Saad}
\date{Source: arXiv:2607.28245v1 (CC BY 4.0)}
\begin{document}
\maketitle

\noindent\emph{Source and licence.} Euler-type recurrence relations for partition functions with congruence conditions. Version arXiv:2607.28245v1 (CC BY 4.0), distributed by arXiv under the Creative Commons Attribution 4.0 International licence, http://creativecommons.org/licenses/by/4.0/, as recorded in the arXiv licence metadata for this exact version and shown on the abstract page https://arxiv.org/abs/2607.28245v1. Full licence notice: \texttt{licenses/arxiv-2607.28245-CC-BY-4.0.txt}.\par\smallskip

\noindent\textit{Editorial scope.} Counted unit: the article's lemma lem:cuspParamLemma (source0.tex lines 923--929). The article's complete original proof of it is delivered with it (source0.tex lines 931--997, one \texttt{proof} environment of 67 lines). No mathematics is written, repaired or re-derived: the statement and the proof are the article's own lines, in the article's own notation, and no retained line is substituted or re-flowed.  Retained uncounted as the source-local context the proof needs: lines 595--597.  Nothing was omitted from any retained span, so the package carries no omission record.  No retained line contains a citation, so the wrapper carries no bibliography and no bibliography entry.  No retained line contains a figure environment, an image-inclusion command or drawing code, so no image is delivered and the package declares no asset dependency.  Editorial formatting only, in addition: an AMS \texttt{amsart} preamble that restates the article's own declarations and macros used by the retained lines, namely the environments \texttt{lemma} on separate counters, together with the standard packages those lines need.  The retained environments print in this document as Lemma 1 in order of appearance, whereas the article prints the same items with the numbering of its own full document.  The complete native source of the article as distributed in the arXiv e-print for this exact version is bundled unchanged as \texttt{sources/206400/source0.tex}.
		\begin{equation}\label{eq:EtaGenTransAppFourier}
		\varepsilon_{\delta,g}(L_\rho^{-1})\cdot \left(\frac{1}{\eta_{\delta,g}^\ast}\Bigg|_{-1/2}{L_\rho^{-1}}\right)(\tau) = \alpha_0(-f_\rho g)^{-1}\cdot q^{-\frac{P_2\left(\frac{h_\rho g}{\delta}\right)}{2}-\frac{1}{24}}\cdot\sum\limits_{0\leq m<\frac{\delta}{t_\rho}K_\rho} pp_\rho(m)q^{\frac{m}{\delta}} + O(1).
		\end{equation}

	\begin{lemma}\label{lem:cuspParamLemma}
		If $m\ge 1$, $n\in \mathbb Z$, $pp_\rho(0)\neq 0,$  and $\rho=\frac{-h_\rho}{g_\rho}$ is a cusp, then
		\[
		\delta\left(\frac{-m+\widetilde{\kappa_\rho}}{t_\rho}+
		\frac{(\delta-2h_\rho g-2\delta n)^2}{8\delta^2}\right)\in \mathbb Z.
		\]
	\end{lemma}

	\begin{proof}
		Let $\alpha_\rho\in\{0,1,\dots,\delta-1\}$ be the representative of $h_\rho g \pmod{\delta}$.
		By (\ref{eq:EtaGenTransAppFourier}), the expansion of $\frac{1}{\eta_{\delta,g}^\ast}$ at the cusp $\rho$ is
		\[
		\left.\frac{1}{\eta^*_{\delta,g}}\right|_{L_\rho^{-1}}(\tau)
		\doteq
		q^{-\frac1{24}-\frac12P_2(\frac{\alpha_\rho}{\delta})}
		\sum_{u} pp_\rho(u) q^{\frac{u}{\delta}} + O(1).
		\]
		Since the cusp width is $t_\rho$ and the cusp parameter with respect to $\overline{\varepsilon_{\delta,g}}$ is $\widetilde{\kappa_\rho}$, all exponents occurring in this expansion lie in
		\[
		\frac{\widetilde{\kappa_\rho}}{t_\rho}+\frac{1}{t_\rho}\mathbb Z.
		\]
		In particular,
		\[
		\frac{\widetilde{\kappa_\rho}}{t_\rho}\equiv
		-\frac1{24}-\frac12P_2\!\left(\frac{\alpha_\rho}{\delta}\right)
		\pmod{\frac1{t_\rho}}.
		\]
		Multiplying by $\delta$ and using $t_\rho\mid \delta$, we obtain
		\[
		\delta\frac{\widetilde{\kappa_\rho}}{t_\rho}\equiv
		-\frac{\delta}{24}-\frac{\delta}{2}P_2\!\left(\frac{\alpha_\rho}{\delta}\right)
		\pmod{\mathbb Z}.
		\]
		Now
		\[
		P_2\!\left(\frac{\alpha_\rho}{\delta}\right)
		=\frac{\alpha_\rho^2}{\delta^2}-\frac{\alpha_\rho}{\delta}+\frac16,
		\]
		so
		\[
		-\frac{\delta}{24}-\frac{\delta}{2}P_2\!\left(\frac{\alpha_\rho}{\delta}\right)
		=
		-\frac{\alpha_\rho^2}{2\delta}+\frac{\alpha_\rho}{2}-\frac{\delta}{8}.
		\]
		Hence
		\[
		\delta\frac{\widetilde{\kappa_\rho}}{t_\rho}
		+\frac{\alpha_\rho^2}{2\delta}-\frac{\alpha_\rho}{2}+\frac{\delta}{8}\in \mathbb Z.
		\]
		Since $\alpha_\rho\equiv h_\rho g \pmod{\delta}$, one checks directly that
		\[
		\frac{\alpha_\rho^2}{2\delta}-\frac{\alpha_\rho}{2}
		\equiv
		\frac{h_\rho^2g^2}{2\delta}-\frac{h_\rho g}{2}
		\pmod{\mathbb Z},
		\]
		and therefore
		\[
		\delta\frac{\widetilde{\kappa_\rho}}{t_\rho}
		+\frac{h_\rho^2g^2}{2\delta}-\frac{h_\rho g}{2}+\frac{\delta}{8}\in \mathbb Z.
		\]
		Finally,
		\[
		\delta\cdot \frac{(\delta-2h_\rho g-2\delta n)^2}{8\delta^2}
		=
		\frac{\delta}{8}-\frac{h_\rho g}{2}+\frac{h_\rho^2g^2}{2\delta}
		+\bigg(\frac{\delta(n^2-n)}{2}+h_\rho g\,n\bigg),
		\]
		and the term in parentheses is an integer. Thus
		\[
		\delta\left(\frac{-m+\widetilde{\kappa_\rho}}{t_\rho}+
		\frac{(\delta-2h_\rho g-2\delta n)^2}{8\delta^2}\right)\in \mathbb Z,
		\]
		as claimed.
	\end{proof}

\end{document}
